\subsection{The Radical of an Artinian Ring}

PROPOSITION 1. --- Let A be a left Artinian ring. The radical r of A is the largest nilpotent two-sided ideal of A, and the ring A/r is semisimple.

We know (VIII, p.~156, Theorem 1) that every nilpotent two-sided ideal of A is contained in r. Let us prove that r is nilpotent. Since the ring A is left Artinian and the sequence of two-sided ideals \(r^{n}\) is decreasing, there exists an integer \(p \geqslant 0\) such that we have \(r^{p} = r^{p+1}\) . Since A is Artinian, the A-module \(A_{s}\) has finite length (VIII, p.~6, Theorem 1) and is therefore Noetherian. The left ideal \(r^{p}\) is therefore finitely generated. By Nakayama's lemma (Theorem 2 of VIII, p.~158), it follows that \(r^{p} = 0\) .

The ring A/r is without radical (VIII, p.~155, Proposition 5); since it is left Artinian, it is semisimple (VIII, p.~154, Proposition 4).

COROLLARY 1. --- A ring is semisimple if and only if it is left Artinian and has no nilpotent two-sided ideals other than 0.

A semisimple ring is left Artinian and without radical (VIII, p.~154, Proposition 4); it therefore has no nilpotent two-sided ideals other than 0. The converse follows from Proposition 1.

COROLLARY 2. --- In a left Artinian ring A, the radical consists of the elements x such that ax is nilpotent for every a in A.

If A is left Artinian, then the radical is a nilpotent two-sided ideal (Proposition 1). Every left nil ideal is contained in the radical by Jacobson's theorem (VIII, p.~156, Theorem 1). Corollary 2 follows.

